% universal_c0.tex — §3

\section{The Universal Mode}
\label{sec:universal_c0}

\claim{c0_universal}{All four distance probes define the same leading SVD direction $V_0$, whose data amplitude is \czero (cross-probe inner products $> 0.996$).}
\claim{c0_is_omegamh2}{\czero is a near-perfect proxy for \omh ($R^2 > 0.99$).}

The SVD of whitened chain predictions reveals that one mode dominates for every probe. Table~\ref{tab:sigma_c} summarizes the per-mode standard deviations $\sigma_{c_\alpha}$ (Eq.~\ref{eq:sigma_c}).

\begin{deluxetable}{lcccc}
\tablecaption{Leading Two SVD Mode Standard Deviations $\sigma_{c_\alpha}$ by Probe \label{tab:sigma_c}}
\tablehead{
\colhead{Mode} & \colhead{BAO} & \colhead{Union3} & \colhead{Pantheon+} & \colhead{DES-Dovekie}
}
\startdata
$\sigma_{c_0}$ & 1.56 & 0.28 & 0.30 & 0.33 \\
$\sigma_{c_1}$ & 0.063 & 0.0010 & 0.0011 & 0.0009 \\
$\sigma_{c_0}/\sigma_{c_1}$ & 25 & 273 & 275 & 355 \\
\enddata
\end{deluxetable}

Within \LCDM, only the amplitude of the leading mode is measurable for any of the four probes. BAO has $\sigma_{c_0} = 1.56 > 1$: the CMB chain populates the $V_0$ direction by more than one BAO measurement error, so projecting the BAO data onto $V_0$ constrains \omh more tightly than the CMB itself. The three SN datasets have $\sigma_{c_0} \approx 0.28$--$0.33 < 1$: the CMB already constrains $V_0$ tighter than the SN data can --- projecting the SN data onto $V_0$ adds little information about \omh beyond the CMB prior. No probe has a second measurable \LCDM direction ($\sigma_{c_1} \ll 1$ for all, including BAO at $\sigma_{c_1} = 0.06$).

\subsection{Universality and Physical Interpretation}
\label{sec:c0_interpretation}

All four probes define almost the same \czero direction in parameter space. Table~\ref{tab:beta_vectors} shows the $\beta$-vector coefficients (Eq.~\ref{eq:beta_vectors}) for each probe:

\begin{deluxetable}{lccc}
\tablecaption{$\beta$-Vector Coefficients for \czero. In all cases the three-parameter fit achieves $R^2 > 0.999$. \label{tab:beta_vectors}}
\tablehead{
\colhead{Probe} & \colhead{$\beta_{\omh}$} & \colhead{$\beta_{\obh}$} & \colhead{$\beta_{\thetastar}$}
}
\startdata
BAO (DESI) & $-0.90$ & $+0.13$ & $+0.17$ \\
Union3 & $-0.90$ & $+0.18$ & $+0.12$ \\
Pantheon+ & $-0.90$ & $+0.18$ & $+0.12$ \\
DES-Dovekie & $-0.90$ & $+0.18$ & $+0.12$ \\
\enddata
\end{deluxetable}

\evidence{c0_universal}{ev_inner_products}{Inner product matrix: minimum = 0.9969, mean = 0.9984 across all probe pairs.}

To compare what the four probes' leading directions $V_0$ have in common, we must first express them in a common space. Each probe's $V_0$ is a unit vector of length $D_\mathrm{probe}$ in its own whitened observable space (e.g., 13-D for BAO, 22-D for Union3). We map each $V_0$ into a common 3-D CMB parameter space via the $\beta$-vector regression (Eq.~\ref{eq:beta_vectors}): the unit-norm vector $\hat\beta = (\beta_{\Omega_m h^2}, \beta_{\Omega_b h^2}, \beta_{\theta_\star})/\|\beta\|$ identifies the linear combination of CMB parameters that each probe's $V_0$ is sensitive to. The pairwise dot product $\hat\beta^{(A)} \cdot \hat\beta^{(B)}$ is what we call the cross-probe inner product. A value of 1 means the two probes' $V_0$ directions correspond to identical parameter combinations; 0 would mean orthogonal. The cross-probe inner product for \czero exceeds 0.996 for every pair. The three SN datasets define virtually identical \czero directions; BAO differs slightly due to its sensitivity to the sound horizon \rs, but the BAO--SN inner product still exceeds 0.996. This universality is not guaranteed a priori: BAO measures $D/\rs$ ratios while SN measures distance moduli $\mu(z)$ --- at low redshift these depend on different combinations of cosmological parameters. The convergence to the same direction is a consequence of the CMB \LCDM posterior: once the CMB pins down the parameters the residual freedom in low-redshift distance predictions runs almost entirely along one direction, which we loosely call the $\Omega_m h^2$ direction, regardless of which probe is doing the measuring.

The dominance of \omh indicates that \czero is primarily an \omh proxy. Here $R^2$ is the fraction of the chain-sample variance of $c_0$ explained by a linear regression on the listed CMB parameters; $R^2 = 1$ would mean $c_0$ is exactly a linear combination of those parameters across the CMB posterior. Sequential multi-parameter fits quantify this:
\begin{itemize}
\item \omh alone: $R^2 = 0.95$
\item $\omh + \obh$: $R^2 = 0.974$
\item $\omh + \obh + \thetastar$: $R^2 = 0.99998$
\end{itemize}
The first line is from a practical perspective the non-trivial content: \omh alone explains 95\% of the chain variance in \czero. Adding \obh and \thetastar absorbs the residual 5\%, giving $R^2 = 0.99998$. The three-parameter $R^2 \approx 1$ is expected --- the remaining \LCDM parameters (optical depth, primordial amplitude and tilt) do not enter low-redshift distances at the linear level --- but the 0.95 from \omh alone is what makes \czero usable as an \omh proxy. For brevity we will refer to \czero as the \omh mode throughout, with the understanding that it is the specific linear combination of (\omh, \obh, \thetastar) along the CMB degeneracy direction.

\evidence{c0_is_omegamh2}{ev_beta_regression}{Sequential fit: \omh alone $R^2=0.95$, $+\obh$ $R^2=0.974$, $+\thetastar$ $R^2=0.99998$. Beta regression on all four probes.}

The sigma-normalized formula for BAO is:
\begin{multline}
\frac{\czero - \langle \czero \rangle}{\sigma_{\czero}} = -0.90\, \frac{\omh - \langle \omh \rangle}{\sigma_{\omh}} \\
+ 0.13\, \frac{\obh - \langle \obh \rangle}{\sigma_{\obh}} + 0.17\, \frac{\thetastar - \langle \thetastar \rangle}{\sigma_{\thetastar}}
\label{eq:c0_formula}
\end{multline}
The SN $\beta$-vectors are nearly identical (Table~\ref{tab:beta_vectors}), so the sigma-normalized \czero is the same linear combination of cosmological parameters for every probe. 

The difference between probes lies only in the normalization: $\czero = 1$ corresponds to $1\sigma$ of measurement error for that probe, and since BAO and SN have very different measurement precisions, a unit of BAO \czero and a unit of SN \czero correspond to different shifts in \omh. Specifically, because the $\beta$-vectors are almost proportional, the \czero values for any two probes are related by:
\begin{equation}
c_0^\mathrm{SN} \approx \frac{\sigma_{c_0}^\mathrm{SN}}{\sigma_{c_0}^\mathrm{BAO}} \, c_0^\mathrm{BAO}
\label{eq:c0_sn_bao}
\end{equation}
For DES-Dovekie, $\sigma_{c_0}^\mathrm{SN}/\sigma_{c_0}^\mathrm{BAO} = 0.33/1.56 \approx 0.21$: a given cosmological shift produces a $\sim 5\times$ smaller \czero value in SN whitened units than in BAO whitened units, reflecting the lower sensitivity of SN to \omh.

The dominance of \omh follows quantitatively from the logarithmic derivatives of low-redshift distance observables combined with the CMB parameter covariance (Appendix~\ref{app:universality}). We define the \emph{per-sigma sensitivity} of an observable to a parameter $p$ as the fractional change in the observable per $1\sigma$ shift in $p$ along the CMB posterior --- i.e., the log derivative times the CMB fractional error. For $\DM(z)/r_d$ and \omh, this is $\sim 5\times$ larger than for \obh, the product of (i) a log derivative $\partial \ln(\DM/r_d)/\partial \ln \omh \approx 0.8$ which is $\sim 3\times$ larger than $|\partial\ln(\DM/r_d)/\partial\ln \obh| \approx 0.3$, and (ii) a CMB fractional error on \omh of $0.8\%$ which is $\sim 1.6\times$ larger than the $0.5\%$ on \obh. The large \thetastar derivative ($g_{\thetastar} \approx -4$) is suppressed by the tiny fractional error ($0.025\%$). 

The correlations between parameters in the CMB posterior further amplify the \omh dominance: when regressing on \omh alone, the correlated parts of \obh and \thetastar are partially captured, so that \omh alone predicts $R^2 = 0.95$. Appendix~\ref{app:universality} predicts the $\beta$-coefficients in Table~\ref{tab:beta_vectors} to within 2\% for \omh and ${\sim}\,30\%$ for \obh from single-redshift proxy observables. The results in the appendix are presented simply to gain some intuition about the results. 

For SN, the $M_B$ marginalization removes the absolute distance scale, leaving the \emph{shape} of the distance--redshift relation. The relevant proxy is the distance ratio $\DM(0.3)/\DM(0.5)$, which has no $r_d$ dependence and no $M_B$ ambiguity. Its predicted $\beta$ is similar to the BAO proxy's (Appendix~\ref{app:universality}, Table~\ref{tab:beta_comparison}), explaining why the SN $\beta$-vectors in Table~\ref{tab:beta_vectors} nearly match the BAO values. The small difference ($\beta_{\obh} = +0.18$ for SN vs $+0.13$ for BAO) traces to the fact that BAO has an additional \obh sensitivity through $r_d$ that the SN distance ratio lacks.

The MEDI framework \citep{Weiner2026_MEDI} provides the analytic counterpart to this result: it identifies $\omega_m r_d^2$ --- rather than \omh alone --- as the parameter combination that acoustic-scale measurements (BAO calibrated by the CMB) are fundamentally sensitive to. Within \LCDM the drag-epoch sound horizon $r_d$ is nearly fixed by early-universe physics, so $\omega_m r_d^2$ is to good approximation proportional to \omh; this is why \czero emerges as essentially an \omh mode, and also why regressing \czero on $\omega_m r_d^2$ across the \LCDM chain is no more informative than regressing on \omh alone --- the chain does not exercise $r_d$ enough to separate the two. The value of the MEDI identification is that it holds analytically, including in extended models where $r_d$ can shift. See the introduction for a fuller discussion of the relation to previous work.

Figure~\ref{fig:c0_families} shows what varying \czero looks like in observable space. The left panel displays BAO $D_V/\rs$ ratios (normalized to the Planck reference cosmology) for five values of \czero spanning $\pm 5$ in BAO whitened units, which corresponds to $\pm 5 / \sigma_{c_0}^\mathrm{BAO} \approx \pm 3.2$ in units of the ACT \LCDM chain spread in $V_0$ (Table~\ref{tab:sigma_c}); the right panel shows the corresponding SN distance modulus residuals (mean-subtracted to marginalize $M_B$). The \czero values in the two panels do not correspond to the same cosmological parameters: the normalization of \czero is set by each probe's measurement covariance (Eq.~\ref{eq:c0_sn_bao}), so $c_0 = 5$ represents a $\sim 5\times$ larger shift in \omh for BAO than for SN.

\begin{figure*}[t]
\centering
\includegraphics[width=\textwidth]{figures/fig_c0_families.pdf}
\caption{\czero family of curves in observed space. \emph{Left:} BAO $D_V/\rs$ ratio (normalized to the Planck reference cosmology) for five values of \czero spanning $\pm 5$ whitened units ($\approx \pm 3\sigma$ of the ACT chain). Black points: DESI DR2. \emph{Right:} SN distance modulus residuals $\Delta\mu$ (mean-subtracted) for the same \czero range. Points: DES-Dovekie. The \czero normalization differs between panels (Eq.~\ref{eq:c0_sn_bao}): BAO \czero $= 5$ corresponds to a $\sim 5\times$ larger \omh shift than SN \czero $= 5$.}
\label{fig:c0_families}
\end{figure*}

\dataref{ev_inner_products}{calculation/scripts/cross\_dataset\_analysis.py}
